
\documentclass{article}
\usepackage{colortbl}
\usepackage{makecell}
\usepackage{multirow}
\usepackage{supertabular}

\begin{document}

\newcounter{utterance}

\centering \large Interaction Transcript for game `matchit\_ascii', experiment `different\_grid', episode 3 with deepseek{-}v4{-}pro.
\vspace{24pt}

{ \footnotesize  \setcounter{utterance}{1}
\setlength{\tabcolsep}{0pt}
\begin{supertabular}{c@{$\;$}|p{.15\linewidth}@{}p{.15\linewidth}p{.15\linewidth}p{.15\linewidth}p{.15\linewidth}p{.15\linewidth}}
    \# & $\;$A & \multicolumn{4}{c}{Game Master} & $\;\:$B\\
    \hline

    \theutterance \stepcounter{utterance}  
    & & \multicolumn{4}{p{0.6\linewidth}}{
        \cellcolor[rgb]{0.9,0.9,0.9}{
            \makecell[{{p{\linewidth}}}]{
                \texttt{\tiny{[P1$\langle$GM]}}
                \texttt{Bierzesz udział w grze polegającej na wspólnym zgadywaniu.} \\
\texttt{Celem jest ustalenie, czy ta siatka i inny obrazek, który widzę tylko ja, są takie same.} \\
\texttt{Proszę najpierw opisać swoją siatkę. Następnie ja przedstawię swój opis i będziemy mogli zadawać sobie nawzajem pytania dotyczące siatek, aby ustalić, czy są takie same. Teraz zacznij swój krótki opis siatki od "OPIS:", a następnie podaj opis. Nie dodawaj nic więcej. To jest twoja siatka:} \\
\texttt{▢ X X X ▢} \\
\texttt{X ▢ X ▢ X} \\
\texttt{X X ▢ X X} \\
\texttt{X X ▢ X X} \\
\texttt{X X ▢ X X} \\
            }
        }
    }
    & \\ \\

    \theutterance \stepcounter{utterance}  
    & \multicolumn{4}{p{0.6\linewidth}}{
        \cellcolor[rgb]{0.9,0.9,0.9}{
            \makecell[{{p{\linewidth}}}]{
                \texttt{\tiny{[P1$\rangle$GM]}}
                \texttt{OPIS: Siatka 5x5. Pierwszy rząd: pusta, X, X, X, pusta. Drugi rząd: X, pusta, X, pusta, X. Trzeci rząd: X, X, pusta, X, X. Czwarty rząd: X, X, pusta, X, X. Piąty rząd: X, X, pusta, X, X.} \\
            }
        }
    }
    & & \\ \\

    \theutterance \stepcounter{utterance}  
    & & & \multicolumn{2}{p{0.3\linewidth}}{
        \cellcolor[rgb]{0.9,0.9,0.9}{
            \makecell[{{p{\linewidth}}}]{
                \texttt{\tiny{[GM$|$GM]}}
                \texttt{continue} \\
            }
        }
    }
    & & \\ \\

    \theutterance \stepcounter{utterance}  
    & & \multicolumn{4}{p{0.6\linewidth}}{
        \cellcolor[rgb]{0.9,0.9,0.9}{
            \makecell[{{p{\linewidth}}}]{
                \texttt{\tiny{[GM$\rangle$P2]}}
                \texttt{Bierzesz udział w grze polegającej na wspólnym zgadywaniu.} \\
\texttt{Celem jest ustalenie, czy ta siatka i inny obrazek, który widzę tylko ja, są takie same.} \\
\texttt{Proszę najpierw opisać swoją siatkę. Następnie ja przedstawię swój opis i będziemy mogli zadawać sobie nawzajem pytania dotyczące siatek, aby ustalić, czy są takie same. Teraz zacznij swój krótki opis siatki od "OPIS:", a następnie podaj opis. Nie dodawaj nic więcej. To jest twoja siatka:} \\
\texttt{▢ ▢ ▢ ▢ X} \\
\texttt{▢ ▢ ▢ X X} \\
\texttt{▢ ▢ X X X} \\
\texttt{▢ X X X X} \\
\texttt{X X X X X} \\
            }
        }
    }
    & \\ \\

    \theutterance \stepcounter{utterance}  
    & & & \multicolumn{4}{p{0.6\linewidth}}{
        \cellcolor[rgb]{0.9,0.9,0.9}{
            \makecell[{{p{\linewidth}}}]{
                \texttt{\tiny{[GM$\langle$P2]}}
                \texttt{OPIS: Siatka o wymiarach 5 na 5. W lewym górnym rogu znajduje się kwadrat 4x4 pustych pól (▢). Po przekątnej od prawego górnego rogu do lewego dolnego rogu ciągną się pola wypełnione (X), tworząc trójkątny kształt, który rozszerza się w dół i w lewo. Pierwszy wiersz ma 4 puste i 1 X na końcu, drugi ma 3 puste i 2 X, trzeci ma 2 puste i 3 X, czwarty ma 1 puste i 4 X, a piąty wiersz to same X.} \\
            }
        }
    }
     \\ \\

    \theutterance \stepcounter{utterance}  
    & & & \multicolumn{2}{p{0.3\linewidth}}{
        \cellcolor[rgb]{0.9,0.9,0.9}{
            \makecell[{{p{\linewidth}}}]{
                \texttt{\tiny{[GM$|$GM]}}
                \texttt{continue} \\
            }
        }
    }
    & & \\ \\

    \theutterance \stepcounter{utterance}  
    & & \multicolumn{4}{p{0.6\linewidth}}{
        \cellcolor[rgb]{0.9,0.9,0.9}{
            \makecell[{{p{\linewidth}}}]{
                \texttt{\tiny{[GM$\rangle$P2]}}
                \texttt{To jest mój OPISOPIS: Siatka 5x5. Pierwszy rząd: pusta, X, X, X, pusta. Drugi rząd: X, pusta, X, pusta, X. Trzeci rząd: X, X, pusta, X, X. Czwarty rząd: X, X, pusta, X, X. Piąty rząd: X, X, pusta, X, X.} \\
\texttt{Teraz zadaj pytanie, aby dowiedzieć się o nowych aspektach mojej siatki, które mogą różnić się od twojej siatki. Zacznij od "PYTANIE:" i nie dodawaj nic więcej.} \\
            }
        }
    }
    & \\ \\

    \theutterance \stepcounter{utterance}  
    & & & \multicolumn{4}{p{0.6\linewidth}}{
        \cellcolor[rgb]{0.9,0.9,0.9}{
            \makecell[{{p{\linewidth}}}]{
                \texttt{\tiny{[GM$\langle$P2]}}
                \texttt{PYTANIE: Czy w twojej siatce wypełnione pola (X) tworzą jeden ciągły, symetryczny wzór krzyża z pustym środkiem, czy może puste pola są rozproszone w inny sposób w dolnych rzędach?} \\
            }
        }
    }
     \\ \\

    \theutterance \stepcounter{utterance}  
    & & & \multicolumn{2}{p{0.3\linewidth}}{
        \cellcolor[rgb]{0.9,0.9,0.9}{
            \makecell[{{p{\linewidth}}}]{
                \texttt{\tiny{[GM$|$GM]}}
                \texttt{continue} \\
            }
        }
    }
    & & \\ \\

    \theutterance \stepcounter{utterance}  
    & & \multicolumn{4}{p{0.6\linewidth}}{
        \cellcolor[rgb]{0.9,0.9,0.9}{
            \makecell[{{p{\linewidth}}}]{
                \texttt{\tiny{[P1$\langle$GM]}}
                \texttt{To jest mój OPISOPIS: Siatka o wymiarach 5 na 5. W lewym górnym rogu znajduje się kwadrat 4x4 pustych pól (▢). Po przekątnej od prawego górnego rogu do lewego dolnego rogu ciągną się pola wypełnione (X), tworząc trójkątny kształt, który rozszerza się w dół i w lewo. Pierwszy wiersz ma 4 puste i 1 X na końcu, drugi ma 3 puste i 2 X, trzeci ma 2 puste i 3 X, czwarty ma 1 puste i 4 X, a piąty wiersz to same X.} \\
\texttt{PYTANIE: Czy w twojej siatce wypełnione pola (X) tworzą jeden ciągły, symetryczny wzór krzyża z pustym środkiem, czy może puste pola są rozproszone w inny sposób w dolnych rzędach?Rozpocznij odpowiedź od "ODPOWIEDŹ:" i nie dodawaj nic więcej.} \\
            }
        }
    }
    & \\ \\

    \theutterance \stepcounter{utterance}  
    & \multicolumn{4}{p{0.6\linewidth}}{
        \cellcolor[rgb]{0.9,0.9,0.9}{
            \makecell[{{p{\linewidth}}}]{
                \texttt{\tiny{[P1$\rangle$GM]}}
                \texttt{ODPOWIEDŹ: Puste pola nie są rozproszone w dolnych rzędach. Wypełnione pola (X) tworzą jeden ciągły, symetryczny wzór krzyża z pustym środkiem. W dolnych trzech rzędach (rząd 3, 4 i 5) puste pole znajduje się dokładnie w środkowej kolumnie, a cała reszta to X.} \\
            }
        }
    }
    & & \\ \\

    \theutterance \stepcounter{utterance}  
    & & & \multicolumn{2}{p{0.3\linewidth}}{
        \cellcolor[rgb]{0.9,0.9,0.9}{
            \makecell[{{p{\linewidth}}}]{
                \texttt{\tiny{[GM$|$GM]}}
                \texttt{continue} \\
            }
        }
    }
    & & \\ \\

    \theutterance \stepcounter{utterance}  
    & & \multicolumn{4}{p{0.6\linewidth}}{
        \cellcolor[rgb]{0.9,0.9,0.9}{
            \makecell[{{p{\linewidth}}}]{
                \texttt{\tiny{[P1$\langle$GM]}}
                \texttt{Teraz zadaj pytanie, aby dowiedzieć się o nowych aspektach mojej siatki, które mogą różnić się od twojej siatki. Zacznij od "PYTANIE:" i nie dodawaj nic więcej.} \\
            }
        }
    }
    & \\ \\

    \theutterance \stepcounter{utterance}  
    & \multicolumn{4}{p{0.6\linewidth}}{
        \cellcolor[rgb]{0.9,0.9,0.9}{
            \makecell[{{p{\linewidth}}}]{
                \texttt{\tiny{[P1$\rangle$GM]}}
                \texttt{PYTANIE: Czy w twojej siatce, w pierwszym rzędzie od góry, puste pola znajdują się dokładnie na pierwszej i ostatniej pozycji (po lewej i po prawej stronie), a trzy środkowe pozycje są wypełnione (X)?} \\
            }
        }
    }
    & & \\ \\

    \theutterance \stepcounter{utterance}  
    & & & \multicolumn{2}{p{0.3\linewidth}}{
        \cellcolor[rgb]{0.9,0.9,0.9}{
            \makecell[{{p{\linewidth}}}]{
                \texttt{\tiny{[GM$|$GM]}}
                \texttt{continue} \\
            }
        }
    }
    & & \\ \\

    \theutterance \stepcounter{utterance}  
    & & \multicolumn{4}{p{0.6\linewidth}}{
        \cellcolor[rgb]{0.9,0.9,0.9}{
            \makecell[{{p{\linewidth}}}]{
                \texttt{\tiny{[GM$\rangle$P2]}}
                \texttt{ODPOWIEDŹ: Puste pola nie są rozproszone w dolnych rzędach. Wypełnione pola (X) tworzą jeden ciągły, symetryczny wzór krzyża z pustym środkiem. W dolnych trzech rzędach (rząd 3, 4 i 5) puste pole znajduje się dokładnie w środkowej kolumnie, a cała reszta to X.} \\
\texttt{PYTANIE: Czy w twojej siatce, w pierwszym rzędzie od góry, puste pola znajdują się dokładnie na pierwszej i ostatniej pozycji (po lewej i po prawej stronie), a trzy środkowe pozycje są wypełnione (X)?Rozpocznij odpowiedź od "ODPOWIEDŹ:" i nie dodawaj nic więcej.} \\
            }
        }
    }
    & \\ \\

    \theutterance \stepcounter{utterance}  
    & & & \multicolumn{4}{p{0.6\linewidth}}{
        \cellcolor[rgb]{0.9,0.9,0.9}{
            \makecell[{{p{\linewidth}}}]{
                \texttt{\tiny{[GM$\langle$P2]}}
                \texttt{ODPOWIEDŹ: Nie. W mojej siatce pierwszy rząd od góry ma puste pola na pierwszych czterech pozycjach od lewej, a wypełnione (X) jest tylko ostatnie, piąte pole po prawej stronie.} \\
            }
        }
    }
     \\ \\

    \theutterance \stepcounter{utterance}  
    & & & \multicolumn{2}{p{0.3\linewidth}}{
        \cellcolor[rgb]{0.9,0.9,0.9}{
            \makecell[{{p{\linewidth}}}]{
                \texttt{\tiny{[GM$|$GM]}}
                \texttt{continue} \\
            }
        }
    }
    & & \\ \\

    \theutterance \stepcounter{utterance}  
    & & \multicolumn{4}{p{0.6\linewidth}}{
        \cellcolor[rgb]{0.9,0.9,0.9}{
            \makecell[{{p{\linewidth}}}]{
                \texttt{\tiny{[GM$\rangle$P2]}}
                \texttt{Teraz zadaj pytanie, aby dowiedzieć się o nowych aspektach mojej siatki, które mogą różnić się od twojej siatki. Zacznij od "PYTANIE:" i nie dodawaj nic więcej.} \\
            }
        }
    }
    & \\ \\

    \theutterance \stepcounter{utterance}  
    & & & \multicolumn{4}{p{0.6\linewidth}}{
        \cellcolor[rgb]{0.9,0.9,0.9}{
            \makecell[{{p{\linewidth}}}]{
                \texttt{\tiny{[GM$\langle$P2]}}
                \texttt{PYTANIE: Czy w twojej siatce drugi rząd od góry ma wypełnione pola (X) na pierwszej, trzeciej i piątej pozycji, a puste pola na drugiej i czwartej?} \\
            }
        }
    }
     \\ \\

    \theutterance \stepcounter{utterance}  
    & & & \multicolumn{2}{p{0.3\linewidth}}{
        \cellcolor[rgb]{0.9,0.9,0.9}{
            \makecell[{{p{\linewidth}}}]{
                \texttt{\tiny{[GM$|$GM]}}
                \texttt{continue} \\
            }
        }
    }
    & & \\ \\

    \theutterance \stepcounter{utterance}  
    & & \multicolumn{4}{p{0.6\linewidth}}{
        \cellcolor[rgb]{0.9,0.9,0.9}{
            \makecell[{{p{\linewidth}}}]{
                \texttt{\tiny{[P1$\langle$GM]}}
                \texttt{ODPOWIEDŹ: Nie. W mojej siatce pierwszy rząd od góry ma puste pola na pierwszych czterech pozycjach od lewej, a wypełnione (X) jest tylko ostatnie, piąte pole po prawej stronie.} \\
\texttt{PYTANIE: Czy w twojej siatce drugi rząd od góry ma wypełnione pola (X) na pierwszej, trzeciej i piątej pozycji, a puste pola na drugiej i czwartej?Rozpocznij odpowiedź od "ODPOWIEDŹ:" i nie dodawaj nic więcej.} \\
            }
        }
    }
    & \\ \\

    \theutterance \stepcounter{utterance}  
    & \multicolumn{4}{p{0.6\linewidth}}{
        \cellcolor[rgb]{0.9,0.9,0.9}{
            \makecell[{{p{\linewidth}}}]{
                \texttt{\tiny{[P1$\rangle$GM]}}
                \texttt{ODPOWIEDŹ: Tak, dokładnie tak. W drugim rzędzie od góry X są na pierwszej, trzeciej i piątej pozycji, a puste na drugiej i czwartej.} \\
            }
        }
    }
    & & \\ \\

    \theutterance \stepcounter{utterance}  
    & & & \multicolumn{2}{p{0.3\linewidth}}{
        \cellcolor[rgb]{0.9,0.9,0.9}{
            \makecell[{{p{\linewidth}}}]{
                \texttt{\tiny{[GM$|$GM]}}
                \texttt{continue} \\
            }
        }
    }
    & & \\ \\

    \theutterance \stepcounter{utterance}  
    & & \multicolumn{4}{p{0.6\linewidth}}{
        \cellcolor[rgb]{0.9,0.9,0.9}{
            \makecell[{{p{\linewidth}}}]{
                \texttt{\tiny{[P1$\langle$GM]}}
                \texttt{Teraz zadaj pytanie, aby dowiedzieć się o nowych aspektach mojej siatki, które mogą różnić się od twojej siatki. Zacznij od "PYTANIE:" i nie dodawaj nic więcej.} \\
            }
        }
    }
    & \\ \\

    \theutterance \stepcounter{utterance}  
    & \multicolumn{4}{p{0.6\linewidth}}{
        \cellcolor[rgb]{0.9,0.9,0.9}{
            \makecell[{{p{\linewidth}}}]{
                \texttt{\tiny{[P1$\rangle$GM]}}
                \texttt{PYTANIE: Czy w twojej siatce w trzecim, czwartym i piątym rzędzie od góry puste pola znajdują się wyłącznie na pozycjach trzeciej i czwartej od lewej, czy może występują również na innych pozycjach?} \\
            }
        }
    }
    & & \\ \\

    \theutterance \stepcounter{utterance}  
    & & & \multicolumn{2}{p{0.3\linewidth}}{
        \cellcolor[rgb]{0.9,0.9,0.9}{
            \makecell[{{p{\linewidth}}}]{
                \texttt{\tiny{[GM$|$GM]}}
                \texttt{continue} \\
            }
        }
    }
    & & \\ \\

    \theutterance \stepcounter{utterance}  
    & & \multicolumn{4}{p{0.6\linewidth}}{
        \cellcolor[rgb]{0.9,0.9,0.9}{
            \makecell[{{p{\linewidth}}}]{
                \texttt{\tiny{[GM$\rangle$P2]}}
                \texttt{ODPOWIEDŹ: Tak, dokładnie tak. W drugim rzędzie od góry X są na pierwszej, trzeciej i piątej pozycji, a puste na drugiej i czwartej.} \\
\texttt{PYTANIE: Czy w twojej siatce w trzecim, czwartym i piątym rzędzie od góry puste pola znajdują się wyłącznie na pozycjach trzeciej i czwartej od lewej, czy może występują również na innych pozycjach?Rozpocznij odpowiedź od "ODPOWIEDŹ:" i nie dodawaj nic więcej.} \\
            }
        }
    }
    & \\ \\

    \theutterance \stepcounter{utterance}  
    & & & \multicolumn{4}{p{0.6\linewidth}}{
        \cellcolor[rgb]{0.9,0.9,0.9}{
            \makecell[{{p{\linewidth}}}]{
                \texttt{\tiny{[GM$\langle$P2]}}
                \texttt{ODPOWIEDŹ: Nie. W mojej siatce w trzecim, czwartym i piątym rzędzie od góry nie ma żadnych pustych pól poza jednym wyjątkiem. W trzecim rzędzie puste pola są na pierwszej i drugiej pozycji od lewej, w czwartym rzędzie puste pole jest tylko na pierwszej pozycji od lewej, a w piątym rzędzie wszystkie pola są wypełnione (X).} \\
            }
        }
    }
     \\ \\

    \theutterance \stepcounter{utterance}  
    & & & \multicolumn{2}{p{0.3\linewidth}}{
        \cellcolor[rgb]{0.9,0.9,0.9}{
            \makecell[{{p{\linewidth}}}]{
                \texttt{\tiny{[GM$|$GM]}}
                \texttt{continue} \\
            }
        }
    }
    & & \\ \\

    \theutterance \stepcounter{utterance}  
    & & \multicolumn{4}{p{0.6\linewidth}}{
        \cellcolor[rgb]{0.9,0.9,0.9}{
            \makecell[{{p{\linewidth}}}]{
                \texttt{\tiny{[GM$\rangle$P2]}}
                \texttt{Teraz zadaj pytanie, aby dowiedzieć się o nowych aspektach mojej siatki, które mogą różnić się od twojej siatki. Zacznij od "PYTANIE:" i nie dodawaj nic więcej.} \\
            }
        }
    }
    & \\ \\

    \theutterance \stepcounter{utterance}  
    & & & \multicolumn{4}{p{0.6\linewidth}}{
        \cellcolor[rgb]{0.9,0.9,0.9}{
            \makecell[{{p{\linewidth}}}]{
                \texttt{\tiny{[GM$\langle$P2]}}
                \texttt{PYTANIE: Czy w twojej siatce, jeśli spojrzeć na trzeci rząd od góry, puste pole znajduje się dokładnie w środkowej, trzeciej kolumnie, czy może puste są też inne pozycje poza tą jedną?} \\
            }
        }
    }
     \\ \\

    \theutterance \stepcounter{utterance}  
    & & & \multicolumn{2}{p{0.3\linewidth}}{
        \cellcolor[rgb]{0.9,0.9,0.9}{
            \makecell[{{p{\linewidth}}}]{
                \texttt{\tiny{[GM$|$GM]}}
                \texttt{continue} \\
            }
        }
    }
    & & \\ \\

    \theutterance \stepcounter{utterance}  
    & & \multicolumn{4}{p{0.6\linewidth}}{
        \cellcolor[rgb]{0.9,0.9,0.9}{
            \makecell[{{p{\linewidth}}}]{
                \texttt{\tiny{[P1$\langle$GM]}}
                \texttt{ODPOWIEDŹ: Nie. W mojej siatce w trzecim, czwartym i piątym rzędzie od góry nie ma żadnych pustych pól poza jednym wyjątkiem. W trzecim rzędzie puste pola są na pierwszej i drugiej pozycji od lewej, w czwartym rzędzie puste pole jest tylko na pierwszej pozycji od lewej, a w piątym rzędzie wszystkie pola są wypełnione (X).} \\
\texttt{PYTANIE: Czy w twojej siatce, jeśli spojrzeć na trzeci rząd od góry, puste pole znajduje się dokładnie w środkowej, trzeciej kolumnie, czy może puste są też inne pozycje poza tą jedną?Rozpocznij odpowiedź od "ODPOWIEDŹ:" i nie dodawaj nic więcej.} \\
            }
        }
    }
    & \\ \\

    \theutterance \stepcounter{utterance}  
    & \multicolumn{4}{p{0.6\linewidth}}{
        \cellcolor[rgb]{0.9,0.9,0.9}{
            \makecell[{{p{\linewidth}}}]{
                \texttt{\tiny{[P1$\rangle$GM]}}
                \texttt{ODPOWIEDŹ: W trzecim rzędzie od góry puste pole znajduje się dokładnie w środkowej, trzeciej kolumnie. Wszystkie pozostałe pozycje w tym rzędzie są wypełnione (X).} \\
            }
        }
    }
    & & \\ \\

    \theutterance \stepcounter{utterance}  
    & & & \multicolumn{2}{p{0.3\linewidth}}{
        \cellcolor[rgb]{0.9,0.9,0.9}{
            \makecell[{{p{\linewidth}}}]{
                \texttt{\tiny{[GM$|$GM]}}
                \texttt{continue} \\
            }
        }
    }
    & & \\ \\

    \theutterance \stepcounter{utterance}  
    & & \multicolumn{4}{p{0.6\linewidth}}{
        \cellcolor[rgb]{0.9,0.9,0.9}{
            \makecell[{{p{\linewidth}}}]{
                \texttt{\tiny{[P1$\langle$GM]}}
                \texttt{Teraz zadaj pytanie, aby dowiedzieć się o nowych aspektach mojej siatki, które mogą różnić się od twojej siatki. Zacznij od "PYTANIE:" i nie dodawaj nic więcej.} \\
            }
        }
    }
    & \\ \\

    \theutterance \stepcounter{utterance}  
    & \multicolumn{4}{p{0.6\linewidth}}{
        \cellcolor[rgb]{0.9,0.9,0.9}{
            \makecell[{{p{\linewidth}}}]{
                \texttt{\tiny{[P1$\rangle$GM]}}
                \texttt{PYTANIE: Czy w twojej siatce w czwartym rzędzie od góry puste pola znajdują się na pierwszej i drugiej pozycji od lewej, a pozostałe trzy pozycje są wypełnione (X)?} \\
            }
        }
    }
    & & \\ \\

    \theutterance \stepcounter{utterance}  
    & & & \multicolumn{2}{p{0.3\linewidth}}{
        \cellcolor[rgb]{0.9,0.9,0.9}{
            \makecell[{{p{\linewidth}}}]{
                \texttt{\tiny{[GM$|$GM]}}
                \texttt{continue} \\
            }
        }
    }
    & & \\ \\

    \theutterance \stepcounter{utterance}  
    & & \multicolumn{4}{p{0.6\linewidth}}{
        \cellcolor[rgb]{0.9,0.9,0.9}{
            \makecell[{{p{\linewidth}}}]{
                \texttt{\tiny{[GM$\rangle$P2]}}
                \texttt{ODPOWIEDŹ: W trzecim rzędzie od góry puste pole znajduje się dokładnie w środkowej, trzeciej kolumnie. Wszystkie pozostałe pozycje w tym rzędzie są wypełnione (X).} \\
\texttt{PYTANIE: Czy w twojej siatce w czwartym rzędzie od góry puste pola znajdują się na pierwszej i drugiej pozycji od lewej, a pozostałe trzy pozycje są wypełnione (X)?Rozpocznij odpowiedź od "ODPOWIEDŹ:" i nie dodawaj nic więcej.} \\
            }
        }
    }
    & \\ \\

    \theutterance \stepcounter{utterance}  
    & & & \multicolumn{4}{p{0.6\linewidth}}{
        \cellcolor[rgb]{0.9,0.9,0.9}{
            \makecell[{{p{\linewidth}}}]{
                \texttt{\tiny{[GM$\langle$P2]}}
                \texttt{ODPOWIEDŹ: Nie. W mojej siatce w czwartym rzędzie od góry puste pole znajduje się tylko na pierwszej pozycji od lewej, a pozostałe cztery pozycje (druga, trzecia, czwarta i piąta) są wypełnione (X).} \\
            }
        }
    }
     \\ \\

    \theutterance \stepcounter{utterance}  
    & & & \multicolumn{2}{p{0.3\linewidth}}{
        \cellcolor[rgb]{0.9,0.9,0.9}{
            \makecell[{{p{\linewidth}}}]{
                \texttt{\tiny{[GM$|$GM]}}
                \texttt{continue} \\
            }
        }
    }
    & & \\ \\

    \theutterance \stepcounter{utterance}  
    & & \multicolumn{4}{p{0.6\linewidth}}{
        \cellcolor[rgb]{0.9,0.9,0.9}{
            \makecell[{{p{\linewidth}}}]{
                \texttt{\tiny{[GM$\rangle$P2]}}
                \texttt{Teraz podejmij decyzję. Jak myślisz: czy Twoja siatka i druga opisana siatka są takie same? Napisz "DECYZJA: ta sama siatka" jeśli uważasz, że to ta sama siatka, albo "DECYZJA: różne siatki" jeśli uważasz, że to różne siatki. Nie dodawaj niczego więcej.} \\
            }
        }
    }
    & \\ \\

    \theutterance \stepcounter{utterance}  
    & & & \multicolumn{4}{p{0.6\linewidth}}{
        \cellcolor[rgb]{0.9,0.9,0.9}{
            \makecell[{{p{\linewidth}}}]{
                \texttt{\tiny{[GM$\langle$P2]}}
                \texttt{DECYZJA: różne siatki} \\
            }
        }
    }
     \\ \\

    \theutterance \stepcounter{utterance}  
    & & & \multicolumn{2}{p{0.3\linewidth}}{
        \cellcolor[rgb]{0.9,0.9,0.9}{
            \makecell[{{p{\linewidth}}}]{
                \texttt{\tiny{[GM$|$GM]}}
                \texttt{continue} \\
            }
        }
    }
    & & \\ \\

    \theutterance \stepcounter{utterance}  
    & & & \multicolumn{2}{p{0.3\linewidth}}{
        \cellcolor[rgb]{0.9,0.9,0.9}{
            \makecell[{{p{\linewidth}}}]{
                \texttt{\tiny{[GM$|$GM]}}
                \texttt{success} \\
            }
        }
    }
    & & \\ \\

    \theutterance \stepcounter{utterance}  
    & & \multicolumn{4}{p{0.6\linewidth}}{
        \cellcolor[rgb]{0.9,0.9,0.9}{
            \makecell[{{p{\linewidth}}}]{
                \texttt{\tiny{[P1$\langle$GM]}}
                \texttt{ODPOWIEDŹ: Nie. W mojej siatce w czwartym rzędzie od góry puste pole znajduje się tylko na pierwszej pozycji od lewej, a pozostałe cztery pozycje (druga, trzecia, czwarta i piąta) są wypełnione (X).} \\
\texttt{Teraz podejmij decyzję. Jak myślisz: czy Twoja siatka i druga opisana siatka są takie same? Napisz "DECYZJA: ta sama siatka" jeśli uważasz, że to ta sama siatka, albo "DECYZJA: różne siatki" jeśli uważasz, że to różne siatki. Nie dodawaj niczego więcej.} \\
            }
        }
    }
    & \\ \\

    \theutterance \stepcounter{utterance}  
    & \multicolumn{4}{p{0.6\linewidth}}{
        \cellcolor[rgb]{0.9,0.9,0.9}{
            \makecell[{{p{\linewidth}}}]{
                \texttt{\tiny{[P1$\rangle$GM]}}
                \texttt{DECYZJA: różne siatki} \\
            }
        }
    }
    & & \\ \\

    \theutterance \stepcounter{utterance}  
    & & & \multicolumn{2}{p{0.3\linewidth}}{
        \cellcolor[rgb]{0.9,0.9,0.9}{
            \makecell[{{p{\linewidth}}}]{
                \texttt{\tiny{[GM$|$GM]}}
                \texttt{continue} \\
            }
        }
    }
    & & \\ \\

    \theutterance \stepcounter{utterance}  
    & & & \multicolumn{2}{p{0.3\linewidth}}{
        \cellcolor[rgb]{0.9,0.9,0.9}{
            \makecell[{{p{\linewidth}}}]{
                \texttt{\tiny{[GM$|$GM]}}
                \texttt{success} \\
            }
        }
    }
    & & \\ \\

\end{supertabular}
}

\end{document}
